% ROZDZIAŁ 2 -- fragment do dołączenia do głównego skryptu.
% W preambule dokumentu umieść:
% \usepackage{listings}
% \lstset{
%   language=C++,
%   basicstyle=\ttfamily\small,
%   columns=fullflexible,
%   keepspaces=true,
%   breaklines=true,
%   showstringspaces=false,
%   frame=single,
%   tabsize=4
% }
% Przykładowe programy zakładają C++20, o ile nie podano inaczej.

\chapter{Specyfikatory i kwalifikatory deklaracji w C++}
\label{chap:specyfikatory-cpp}

\section{Wprowadzenie: czym właściwie są \texttt{static}, \texttt{const} i \texttt{inline}?}

Zanim zapamiętamy słowa kluczowe, warto zadać trzy pytania:
czy wartość wolno zmienić, jak długo istnieje obiekt i czy inne pliki
programu widzą ten sam obiekt. \texttt{const} odpowiada głównie na
pierwsze pytanie. \texttt{static} może wpływać na czas życia albo
widoczność nazwy, zależnie od miejsca deklaracji. \texttt{inline}
pozwala umieścić odpowiednią definicję w nagłówku dołączanym przez
kilka plików \texttt{.cpp}.

Słowa kluczowe \texttt{static}, \texttt{extern}, \texttt{inline}, \texttt{const},
\texttt{constexpr} i podobne nie są jednym rodzajem ``modyfikatorów''.
Mogą określać różne, w znacznej mierze niezależne własności deklaracji:
\begin{itemize}
    \item \emph{czas przechowywania} obiektu (\emph{storage duration}), czyli
    jak długo istnieje jego pamięć;
    \item \emph{łączność} nazwy (\emph{linkage}), czyli czy ta sama nazwa
    może oznaczać ten sam byt w różnych jednostkach translacji;
    \item \emph{modyfikowalność} obiektu lub dostępu przez dany typ;
    \item \emph{możliwość obliczenia wyrażenia podczas kompilacji};
    \item \emph{reguły wielokrotnego definiowania} funkcji lub zmiennej;
    \item \emph{własności funkcji i metod}, np. zakaz niejawnej konwersji
    albo gwarancje dotyczące wyjątków.
\end{itemize}

\paragraph{Jednostka translacji.}
W uproszczeniu plik \texttt{.cpp} po rozwinięciu dyrektyw
\texttt{\#include} tworzy jednostkę translacji
(\emph{translation unit}, TU). Plik nagłówkowy bywa włączany do wielu
jednostek translacji, dlatego umieszczanie w nim definicji wymaga uwagi.
Praktycznie: jedna TU to plik \texttt{.cpp} razem z treścią nagłówków,
którą kompilator widzi podczas jego kompilowania.

\paragraph{Czas przechowywania.}
Rozróżniamy w szczególności:
\begin{itemize}
    \item \emph{automatyczny} --- zwykle zmienne lokalne, tworzone przy
    wejściu w odpowiedni zakres i niszczone przy jego opuszczaniu;
    \item \emph{statyczny} --- obiekt istnieje przez czas działania programu;
    \item \emph{wątkowy} --- oddzielny obiekt istnieje dla każdego wątku;
    \item \emph{dynamiczny} --- czas życia kontrolowany jawnie lub przez
    obiekty zarządzające zasobem (np. \texttt{std::unique\_ptr}).
\end{itemize}

\paragraph{Łączność.}
Nazwa może nie mieć łączności (typowo zmienna lokalna), mieć
\emph{łączność wewnętrzną} (\emph{internal linkage}: dostępna jako ten
sam byt tylko w obrębie TU) albo \emph{łączność zewnętrzną}
(\emph{external linkage}: może odnosić się do jednego bytu w wielu TU).
C++20 wprowadza również łączność modułową; w tym rozdziale zajmujemy się
głównie klasycznym modelem plików nagłówkowych i źródłowych.

\section{\texttt{const}: zakaz modyfikacji przez dany dostęp}

\begin{center}\small
\begin{tabularx}{\textwidth}{@{}p{0.30\textwidth}X@{}}
\toprule
Specyfikator & Znaczenie \\
\midrule
\code{const} & Zabrania modyfikacji przez dany dostęp; wartość nie musi być znana podczas kompilacji. \\
\code{constexpr} & Pozwala użyć wartości lub funkcji w obliczeniu stałym, gdy spełnione są jego warunki. \\
\code{consteval} & Wymaga obliczenia wywołania funkcji w czasie kompilacji. \\
\code{constinit} & Wymaga statycznej inicjalizacji zmiennej o statycznym lub wątkowym czasie przechowywania; nie czyni jej stałą. \\
\code{mutable} & Pozwala zmieniać wskazaną składową także przez stały obiekt. \\
\bottomrule
\end{tabularx}
\end{center}

Deklaracja \texttt{const} kwalifikuje typ. Dla obiektu oznacza, że
nie wolno zmieniać jego wartości przez dostęp, w którym typ jest
\texttt{const} (z uwzględnieniem szczególnych zasad np. \texttt{mutable}).
Samo \texttt{const} \emph{nie gwarantuje} obliczenia wartości w czasie kompilacji.

\begin{lstlisting}
int read_value();

void example() {
    const int a = 10;
    // a = 11;              // error: a is const

    const int b = read_value();  // OK: runtime initialization
    int c = read_value();
    c += b;                      // OK: c is mutable
}
\end{lstlisting}

\subsection{Wskaźniki, referencje i pozycja \texttt{const}}

Różnica między kwalifikacją wskaźnika i wskazywanego obiektu jest
fundamentalna:

\begin{lstlisting}
int a = 1;
int b = 2;

const int* p = &a;  // pointer to const int
p = &b;             // OK: pointer can change
// *p = 3;          // error: pointee is const through p

int* const q = &a;  // const pointer to int
*q = 3;             // OK: a can change
// q = &b;          // error: pointer is const

const int* const r = &a;  // both are const through r
const int& ref = a;      // read-only access through ref
\end{lstlisting}

\paragraph{Argumenty i metody.}
Dla dużych obiektów często używamy \texttt{const T\&}, aby uniknąć
kopiowania i zadeklarować brak modyfikacji przez referencję.
Kwalifikator \texttt{const} za listą parametrów metody odnosi się
do obiektu \texttt{*this}:

\begin{lstlisting}
#include <string>

class Person {
    std::string name_;
public:
    explicit Person(std::string name) : name_(name) {}

    const std::string& name() const { return name_; }
    void rename(std::string s) { name_ = s; }
};
\end{lstlisting}

\paragraph{Kiedy stosować?}
Używaj \texttt{const} dla wartości, wskaźników, referencji i metod,
jeśli modyfikacja przez dany interfejs nie powinna być dozwolona.
Nie używaj go jako zamiennika \texttt{constexpr}: zwykła stała
zainicjalizowana wynikiem funkcji może być wyznaczona dopiero podczas
wykonywania programu.

\section{\texttt{constexpr}: stałe i funkcje dostępne dla obliczeń stałych}

\texttt{constexpr} dla zmiennej wymaga inicjalizacji nadającej się
do obliczenia jako wyrażenie stałe; taka zmienna jest kwalifikowana
\texttt{const}. Dla funkcji \texttt{constexpr} oznacza \emph{możliwość}
wywołania w kontekście wymagającym wyniku obliczanego w czasie kompilacji,
a nie obowiązek obliczenia każdego wywołania podczas kompilacji.

\begin{lstlisting}
#include <array>

constexpr int square(int x) {
    return x * x;
}

int read_value();

void example() {
    constexpr int n = square(5);    // constant expression
    std::array<int, n> data{};     // valid array extent

    int x = read_value();
    int y = square(x);             // runtime call is valid

    // constexpr int z = square(x);  // error: x not constant
    (void)data;
    (void)y;
}
\end{lstlisting}

\paragraph{Ważne rozróżnienie.}
Zwykle \texttt{const int n = 5;} także może uczestniczyć w wyrażeniach
stałych, bo spełnia odpowiednie reguły języka; nie oznacza to jednak,
że \emph{każda} zmienna \texttt{const} jest stała czasu kompilacji.
\texttt{constexpr} wyraża ten zamiar jednoznacznie i jest sprawdzany
przez kompilator.

\paragraph{Kiedy stosować?}
Używaj \texttt{constexpr}, gdy wartość ma być wykorzystywana jako stała
czasu kompilacji, np. do rozmiarów tablic \texttt{std::array} i argumentów
szablonów. Nadaje się także do \texttt{static\_assert}. Oznaczaj nim również
funkcje, które naturalnie nadają się zarówno do obliczeń stałych,
jak i zwykłych wywołań.

\subsection{\texttt{if constexpr} (od C++17)}

To osobna konstrukcja sterująca, typowo stosowana w szablonach.
W odróżnieniu od zwykłego \texttt{if} może odrzucić gałąź podczas
instancjacji szablonu:

\begin{lstlisting}
#include <type_traits>

template<class T>
const char* category() {
    if constexpr (std::is_integral_v<T>) {
        return "integral";
    } else {
        return "non-integral";
    }
}
\end{lstlisting}

\section{\texttt{consteval} i \texttt{constinit} (od C++20)}

\subsection{\texttt{consteval}: wywołanie wymagające stałej ewaluacji}

Funkcja \texttt{consteval} wymaga wyniku obliczanego podczas kompilacji.
Dlatego można przekazać jej stałą, ale zwykle nie wartość wczytaną
dopiero podczas działania programu. Dokładne reguły języka obejmują
również wywołania wewnątrz innych funkcji natychmiastowych.

\begin{lstlisting}
consteval int cube(int x) {
    return x * x * x;
}

int read_value();

void example() {
    constexpr int a = cube(3);  // OK
    int b = cube(4);            // OK: call evaluated as constant
    int x = read_value();
    // int c = cube(x);         // error: x not constant
    (void)b;
    (void)x;
}
\end{lstlisting}

\texttt{consteval} wybieraj wtedy, kiedy wywołanie \emph{musi}
być obliczone w kontekście stałym; nie jest to ogólny przepis
na szybsze wykonywanie programu.

\subsection{\texttt{constinit}: gwarancja inicjalizacji statycznej}

\texttt{constinit} dotyczy zmiennych o statycznym lub wątkowym czasie
przechowywania. Wymusza inicjalizację statyczną, a więc wyklucza
inicjalizację dynamiczną. \emph{Nie czyni zmiennej stałą}:

\begin{lstlisting}
constexpr int twice(int x) { return 2 * x; }

constinit int global_count = twice(10);

void update() {
    global_count += 1;        // OK: mutable variable
}

int read_value();
// constinit int bad = read_value();  // error: dynamic init
\end{lstlisting}

Nie wolno łączyć \texttt{constinit} i \texttt{constexpr} w tej samej
deklaracji. \texttt{constinit} ma sens np. dla mutowalnego stanu
globalnego, który musi być inicjalizowany bez inicjalizacji dynamicznej.
Nie rozwiązuje natomiast problemów synchronizacji dostępu wielu wątków.

\section{\texttt{static}: trzy różne zastosowania}

\begin{center}\small
\begin{tabularx}{\textwidth}{@{}p{0.30\textwidth}X@{}}
\toprule
Zapis & Skutek \\
\midrule
\code{static} lokalne & Daje zmiennej lokalnej czas życia do końca programu; inicjalizacja następuje przy pierwszym przejściu przez deklarację. \\
\code{static} składowa & Tworzy składową wspólną dla obiektów klasy. \\
\code{static} w pliku & Nadaje nazwie na poziomie przestrzeni nazw łączność wewnętrzną. \\
\code{extern} & Deklaruje nazwę o łączności zewnętrznej; definicja może być w innym pliku. \\
\code{inline} & Umożliwia zgodne z ODR definicje funkcji lub zmiennej w wielu jednostkach translacji; nie nakazuje optymalizacji wywołania. \\
\bottomrule
\end{tabularx}
\end{center}

Znaczenie \texttt{static} zależy od miejsca deklaracji. Jest to jedno
z najczęstszych źródeł nieporozumień.

\subsection{Nazwa w przestrzeni nazw: łączność wewnętrzna}

W przestrzeni nazw (w tym globalnej) \texttt{static} nadaje nazwie
łączność wewnętrzną. Definicja umieszczona w dwóch plikach
\texttt{.cpp} opisuje wtedy dwa różne byty.

\begin{lstlisting}
// file: first.cpp
static int cache_size = 100;
static int helper(int x) { return x + 1; }

// file: second.cpp
static int cache_size = 200;  // a different variable
\end{lstlisting}

We współczesnym C++ często preferuje się \emph{anonimową przestrzeń nazw}
dla lokalnych szczegółów implementacji:

\begin{lstlisting}
// implementation.cpp
namespace {
    int cache_size = 100;
    int helper(int x) { return x + 1; }
}
\end{lstlisting}

\paragraph{Czego unikać?}
Nie definiuj \texttt{static} zmiennej w nagłówku z zamiarem stworzenia
\emph{jednego} wspólnego obiektu. Każda TU otrzyma swoją instancję.
Jeżeli celem jest wspólny obiekt w nagłówku, rozważ \texttt{inline}
(od C++17), opisane dalej.

Jeśli \texttt{a.cpp} i \texttt{b.cpp} dołączają ten sam nagłówek,
\texttt{static int x = 0;} tworzy dwa oddzielne \texttt{x}.
\texttt{inline int x = 0;} tworzy jeden wspólny obiekt, a
\texttt{extern int x;} jedynie ogłasza wspólny obiekt, którego
jedna definicja musi znaleźć się w pliku \texttt{.cpp}.

\subsection{Zmienna lokalna: czas przechowywania statyczny}

Wewnątrz funkcji \texttt{static} nie nadaje nazwie łączności zewnętrznej
ani wewnętrznej: lokalna nazwa nadal nie ma łączności. Zmienia natomiast
czas przechowywania obiektu. Inicjalizacja lokalnej zmiennej statycznej,
jeśli nie została wykonana wcześniej, następuje przy pierwszym przejściu
sterowania przez jej deklarację.

\begin{lstlisting}
int next_id() {
    static int counter = 0;
    return ++counter;
}

// next_id() returns 1, then 2, then 3, ...
\end{lstlisting}

Od C++11 jednorazowa inicjalizacja lokalnej zmiennej statycznej jest
bezpieczna wobec współbieżności: tylko jeden wątek przeprowadza
inicjalizację, a pozostałe odpowiednio czekają. \emph{Nie wynika z tego},
że późniejsze modyfikacje są bezpieczne. Przy wywołaniach równoległych
powyższe \texttt{++counter} może powodować wyścig danych.

\begin{lstlisting}
#include <atomic>

int next_id_thread_safe() {
    static std::atomic<int> counter{0};
    return counter.fetch_add(1, std::memory_order_relaxed) + 1;
}
\end{lstlisting}

\paragraph{Kiedy stosować?}
Lokalna \texttt{static} jest dobra dla stanu należącego do funkcji
lub obiektu tworzonego leniwie przy pierwszym użyciu. Unikaj jej,
jeśli miałaby ukrywać zależności, utrudniać testy lub tworzyć
współdzielony mutowalny stan bez synchronizacji.

\subsection{Składowa klasy: wspólna dla wszystkich obiektów}

Statyczne pola klasy należą do klasy, a nie do każdej instancji
oddzielnie. Statyczna metoda nie ma wskaźnika \texttt{this}; dlatego
nie może bezpośrednio odwoływać się do niestatycznych pól.

\begin{lstlisting}
#include <cstddef>

struct Config {
    inline static int attempts = 0;          // C++17
    static constexpr std::size_t limit = 256; // inline since C++17

    static void record_attempt() {
        ++attempts;
    }
};

void example() {
    Config::record_attempt();
    auto n = Config::limit;
    (void)n;
}
\end{lstlisting}

Przed C++17 zwykłe statyczne pole niestałe deklarowano w klasie,
a definiowano raz poza nią:

\begin{lstlisting}
// counter.hpp
struct Counter {
    static int value;
};

// counter.cpp
#include "counter.hpp"
int Counter::value = 0;
\end{lstlisting}

Od C++17 \texttt{inline static} pozwala na definicję z inicjalizacją
bezpośrednio w klasie. Pole \texttt{static constexpr} jest od C++17
niejawnie \texttt{inline} i nie wymaga osobnej definicji poza klasą.
W starszych standardach dodatkowa definicja bywała potrzebna przy
\emph{odr-use}, np. przy pobieraniu adresu.

\section{\texttt{extern}: deklaracja bytu zdefiniowanego gdzie indziej}

Typowy wzorzec to deklaracja w pliku nagłówkowym oraz \emph{jedna}
definicja w pliku źródłowym. \texttt{extern} w takiej deklaracji
wskazuje, że nazwa może odnosić się do bytu z łącznością zewnętrzną;
nie alokuje nowego obiektu.

\begin{lstlisting}
// config.hpp
#pragma once
extern int connection_count;
extern const int max_connections;

// config.cpp
#include "config.hpp"
int connection_count = 0;
const int max_connections = 64;

// main.cpp
#include "config.hpp"
void connect() {
    if (connection_count < max_connections)
        ++connection_count;
}
\end{lstlisting}

Ponieważ deklaracja w \texttt{config.cpp} zawiera wcześniejsze
\texttt{extern}, definicja tej zmiennej ma łączność zewnętrzną.
Natomiast samodzielna definicja zmiennej typu \texttt{const int}
na poziomie przestrzeni nazw ma domyślnie łączność wewnętrzną
(w klasycznym kodzie bez modułów).

\paragraph{Uwaga na definicje.}
Zapis \texttt{extern int x = 5;} \emph{jest definicją}, gdy występuje
na poziomie przestrzeni nazw: inicjalizator zmienia znaczenie deklaracji.
Dlatego nie umieszczamy takiego zapisu w zwykłym nagłówku dołączanym
do wielu TU.

\subsection{Dlaczego \texttt{extern const} nie zastąpi \texttt{constexpr}?}

\begin{lstlisting}
// constants.hpp
extern const int array_size;

// constants.cpp
#include "constants.hpp"
const int array_size = 16;

// client.cpp
#include "constants.hpp"
// std::array<int, array_size> a;  // invalid: not constant here
\end{lstlisting}

Inna TU nie widzi inicjalizatora, więc nie może wykorzystywać tej
zmiennej jako znanej stałej szablonowej. Jeśli wartość musi być
stała czasu kompilacji i wspólna w wielu TU, od C++17 zwykle
wybieramy \texttt{inline constexpr} w nagłówku.

\subsection{\texttt{extern "C"}: łączność językowa}

Jest to inne zastosowanie słowa \texttt{extern}: określa
\emph{łączność językową}, np. przy współpracy z biblioteką C.

\begin{lstlisting}
extern "C" int c_library_function(int);
\end{lstlisting}

Nie oznacza ono, że wywołanie funkcji ma semantykę języka C;
przede wszystkim dotyczy zgodności łączenia i typu funkcji.

\section{\texttt{inline}: dozwolone definicje w wielu TU}

Historycznie \texttt{inline} sugerowało wstawienie ciała funkcji
w miejsce wywołania, ale \emph{nie wymusza} takiej optymalizacji.
Najważniejsza współczesna rola polega na odpowiednich regułach
jednej definicji (\emph{One Definition Rule}, ODR).

\subsection{Funkcje \texttt{inline}}

Definicja funkcji \texttt{inline} może znaleźć się w wielu TU,
np. przez dołączenie nagłówka, jeśli spełnione są wymagania ODR,
w szczególności zgodność definicji. Dla funkcji \texttt{inline}
o łączności zewnętrznej wszystkie takie definicje odnoszą się
do jednej funkcji (w tym do tego samego adresu).

\begin{lstlisting}
// math.hpp
#pragma once
inline int add(int a, int b) {
    return a + b;
}
\end{lstlisting}

Gdybyśmy usunęli \texttt{inline} i dołączyli nagłówek do dwóch
\texttt{.cpp}, mielibyśmy niedozwolone wielokrotne definicje zwykłej
funkcji o łączności zewnętrznej.

Funkcje zdefiniowane w obrębie definicji klasy są zazwyczaj niejawnie
\texttt{inline} (z zastrzeżeniem szczególnych reguł modułów C++20).
Funkcje \texttt{constexpr} i \texttt{consteval} też są niejawnie
\texttt{inline} przy pierwszej deklaracji.

\subsection{Zmienne \texttt{inline} (od C++17)}

Zmienna \texttt{inline} zdefiniowana w nagłówku może oznaczać
\emph{jeden wspólny obiekt} w wielu TU:

\begin{lstlisting}
// state.hpp
#pragma once
#include <atomic>

inline std::atomic<int> requests{0};
inline constexpr int buffer_size = 4096;

inline int register_request() {
    return requests.fetch_add(1, std::memory_order_relaxed) + 1;
}
\end{lstlisting}

To praktyczna podstawa bibliotek \emph{header-only}. Na poziomie
przestrzeni nazw zmienna \texttt{inline constexpr} ma zazwyczaj
łączność zewnętrzną. Sama zmienna \texttt{constexpr} ma zaś
domyślnie łączność wewnętrzną.

\paragraph{Ważne: \texttt{static inline} to nie to samo.}
Na poziomie przestrzeni nazw \texttt{static inline} \texttt{int x = 0;}
ma łączność wewnętrzną: każda TU posiada własne \texttt{x}.
Dodanie \texttt{inline} nie zmienia tego w jeden globalny obiekt.

\subsection{Reguła ODR w praktyce}

\begin{itemize}
    \item Zwykła funkcja lub zmienna o łączności zewnętrznej:
    zasadniczo jedna definicja w całym programie.
    \item Funkcja \texttt{inline}, zmienna \texttt{inline}, a także
    odpowiednio zdefiniowane szablony: dozwolone definicje w wielu TU,
    z dodatkowymi wymaganiami co do ich identyczności.
    \item Sama deklaracja w nagłówku jest zwykle bezpieczna;
    problem dotyczy definicji.
\end{itemize}

\section{Łączenie specyfikatorów: najważniejsze przypadki}

Poniższe zapisy nie są wzajemnie zamienne:

\begin{lstlisting}
const int a = 7;                // const; namespace: internal linkage
constexpr int b = 7;            // compile-time const; internal linkage
inline constexpr int c = 7;     // C++17: shared across TUs
static int d = 7;               // namespace: internal linkage
static constexpr int e = 7;     // compile-time; internal linkage
inline int f = 7;               // C++17: one shared object
extern int g;                   // declaration, definition elsewhere
extern const int h;             // declaration, definition elsewhere
constinit int i = 7;            // C++20: static init, mutable
static constinit int j = 7;     // C++20: internal linkage
\end{lstlisting}

\noindent W powyższym przykładzie zakładamy, że wszystkie deklaracje
znajdują się na poziomie przestrzeni nazw, poza modułami nazwanymi.
W rzeczywistym projekcie \texttt{g} i \texttt{h} muszą mieć definicje,
jeśli są odpowiednio używane; nie należy też kopiować wielu definicji
\texttt{i} do nagłówków.

\begin{center}
\small
\renewcommand{\arraystretch}{1.22}
\begin{tabular}{|p{0.24\textwidth}|p{0.31\textwidth}|p{0.34\textwidth}|}
\hline
\textbf{Konstrukcja} & \textbf{Najważniejsza cecha} & \textbf{Typowe zastosowanie} \\
\hline
\texttt{const} & Brak modyfikacji przez dany dostęp & Parametry, referencje, stałe wartości \\
\hline
\texttt{constexpr} & Możliwość użycia wyniku jako wyrażenia stałego & Parametry szablonów, rozmiary, funkcje obliczeniowe \\
\hline
\texttt{consteval} & Wymagana stała ewaluacja wywołania & Walidacja i generowanie stałych przy kompilacji \\
\hline
\texttt{constinit} & Zakaz inicjalizacji dynamicznej & Mutowalny stan o inicjalizacji statycznej \\
\hline
\texttt{static} lokalne & Jeden obiekt o statycznym czasie przechowywania & Leniwy stan funkcji \\
\hline
\texttt{static} globalne & Łączność wewnętrzna & Szczegóły implementacji w \texttt{.cpp} \\
\hline
\texttt{static} w klasie & Składowa wspólna dla instancji & Liczniki, konfiguracja klasy \\
\hline
\texttt{extern} & Deklaracja nazwy odnoszącej się do innej definicji & Podział \texttt{.hpp}/\texttt{.cpp} \\
\hline
\texttt{inline} & Szczególne reguły wielokrotnych definicji & Funkcje i zmienne w nagłówkach \\
\hline
\end{tabular}
\end{center}

\section{\texttt{thread\_local}: osobny obiekt dla każdego wątku}

\texttt{thread\_local} oznacza wątkowy czas przechowywania. Każdy
wątek posiada własną instancję obiektu, więc modyfikacje tej instancji
nie są współdzielone z innymi wątkami.

\begin{lstlisting}
int next_local_id() {
    thread_local int id = 0;
    return ++id;
}
\end{lstlisting}

Różnica względem \texttt{static int id = 0;} jest istotna:
\texttt{static} daje jedną instancję współdzieloną przez wątki,
\texttt{thread\_local} --- po jednej na wątek.

Na poziomie przestrzeni nazw można użyć \texttt{static thread\_local}
(dla łączności wewnętrznej) albo \texttt{extern thread\_local}
(w deklaracji obiektu zdefiniowanego w innym TU).

\paragraph{Kiedy stosować?}
Dla stanu rzeczywiście związanego z wątkiem, np. lokalnego generatora
lub bufora roboczego. Nie jest to ogólny zamiennik blokad i nie nadaje
się do danych, które muszą być wspólne między wątkami.

\section{\texttt{mutable} i \texttt{volatile}: dwa inne rodzaje kwalifikacji}

\subsection{\texttt{mutable}: kontrolowana modyfikacja wewnątrz obiektu stałego}

\texttt{mutable} dotyczy niestatycznych pól klasy. Pozwala modyfikować
takie pole nawet w metodzie \texttt{const}. Typowy przypadek to blokada,
która jest szczegółem technicznym i nie stanowi logicznego stanu obiektu:

\begin{lstlisting}
#include <mutex>

class SafeCounter {
    mutable std::mutex mutex_;
    int value_ = 0;

public:
    void increment() {
        std::lock_guard<std::mutex> lock(mutex_);
        ++value_;
    }

    int value() const {
        std::lock_guard<std::mutex> lock(mutex_);
        return value_;
    }
};
\end{lstlisting}

Nie stosuj \texttt{mutable} tylko po to, aby obchodzić prawidłowo
zaprojektowany interfejs \texttt{const}. Sam \texttt{mutable}
nie zapewnia bezpieczeństwa wątkowego.

\subsection{\texttt{volatile}: nie jest mechanizmem synchronizacji}

\texttt{volatile} służy do określonych sytuacji, w których dostępy
do obiektu mają być obserwowalnymi operacjami (np. w pewnych
zastosowaniach niskopoziomowych). Nie zapewnia atomowości,
porządkowania pamięci ani braku wyścigów danych w C++.

\begin{lstlisting}
// Not a thread-safe counter:
volatile int wrong_counter = 0;

// For concurrent increments, use e.g.:
#include <atomic>
std::atomic<int> correct_counter{0};
\end{lstlisting}

Dla synchronizacji wątków stosuj \texttt{std::atomic}, muteksy i inne
narzędzia współbieżności, a nie \texttt{volatile}.

\section{\texttt{auto} i \texttt{decltype(auto)}: dedukcja typu}

We współczesnym C++ \texttt{auto} w deklaracji zmiennej z inicjalizatorem
włącza dedukcję typu. Dedukcja przez wartość odrzuca
najwyższy poziom \texttt{const} oraz referencyjność wyrażenia.

Natomiast \texttt{decltype(auto)} zachowuje semantykę \texttt{decltype}.

\begin{lstlisting}
int& shared_value() {
    static int n = 10;
    return n;
}

void example() {
    const int k = 4;
    auto a = k;             // int (copy)
    const auto b = k;       // const int

    auto x = shared_value();           // int: kopia
    auto& original = shared_value();   // int&: ten sam obiekt
    x = 15;                            // nie zmienia shared_value()
    original = 20;                     // zmienia shared_value()
    decltype(auto) y = shared_value(); // int&: także referencja
    y = 21;
    (void)a;
    (void)b;
    (void)x;
}
\end{lstlisting}

Dedukcja typu z \texttt{auto} poprawia czytelność, zwłaszcza
przy iteratorach i rozbudowanych typach. Nie należy jednak
niezauważenie tracić referencji, gdy kod wymaga operowania
na oryginalnym obiekcie. Wtedy używaj np. \texttt{auto\&} lub
\texttt{const auto\&}.

\section{Powiązane specyfikatory metod i funkcji}

Nie są to specyfikatory czasu przechowywania, lecz często występują
obok poprzednich słów kluczowych:

\begin{lstlisting}
#include <cstddef>

struct Base {
    virtual ~Base() = default;
    virtual int evaluate() const noexcept { return 1; }
};

struct Derived final : Base {
    int evaluate() const noexcept override { return 2; }
};

struct Index {
    explicit Index(std::size_t x) : value(x) {}
    std::size_t value;
};

[[nodiscard]] int compute_result() { return 42; }
\end{lstlisting}

\begin{itemize}
    \item \texttt{explicit} kontroluje niejawne konwersje
    przez konstruktory i operatory konwersji (szczegóły poniżej);
    \item \texttt{virtual} umożliwia dynamiczny polimorfizm,
    \texttt{override} pozwala sprawdzić, czy metoda rzeczywiście
    nadpisuje metodę bazową, a \texttt{final} zabrania dalszego
    nadpisania metody lub dziedziczenia po klasie;
    \item \texttt{noexcept} deklaruje funkcję nierzucającą wyjątków
    (naruszenie tej gwarancji prowadzi do \texttt{std::terminate});
    \item \texttt{[[nodiscard]]} pozwala kompilatorowi ostrzegać,
    gdy wynik został zignorowany.

    Atrybut \texttt{[[maybe\_unused]]} pozwala ograniczyć
    ostrzeżenia o nieużywanych bytach.
\end{itemize}

\subsection{\texttt{explicit}: kontrola niejawnych konwersji}

Konstruktor bez \texttt{explicit} może zostać użyty przez kompilator do
niejawnej konwersji, jeżeli pozostałe reguły konwersji na to pozwalają.
\texttt{explicit} wymaga, by w takim przypadku zamiar utworzenia obiektu
był zapisany jawnie. Dotyczy to nie tylko konstruktorów z jednym
parametrem: także konstruktor z kilkoma parametrami może uczestniczyć
w inicjalizacji przez listę.

\begin{center}\small
\begin{tabularx}{\textwidth}{@{}p{0.37\textwidth}X@{}}
\toprule
Zapis & Skutek dla konstruktora \code{explicit Meter(int)} \\
\midrule
\code{Meter a(3)} & Inicjalizacja bezpośrednia: dozwolona. \\
\code{Meter a{3}} & Inicjalizacja bezpośrednia przez listę: dozwolona. \\
\code{Meter a = 3} & Inicjalizacja kopiująca wymagałaby niejawnej konwersji: błąd kompilacji. \\
\code{Meter a = {3}} & Inicjalizacja kopiująca przez listę nie może zakończyć się wyborem konstruktora \code{explicit}: błąd kompilacji. \\
\code{use(Meter{3})} & Jawne utworzenie argumentu: dozwolone. Samo \code{use(3)} byłoby niedozwolone. \\
\bottomrule
\end{tabularx}
\end{center}

\begin{lstlisting}[caption={Konstruktor i operator konwersji explicit (kompletny program, C++20).}]
#include <iostream>

struct Meter {
    explicit Meter(int n) : value(n) {}
    int value;
};

void use(Meter m) {
    std::cout << m.value << '\n';
}

struct Ready {
    bool value;
    explicit operator bool() const { return value; }
};

int main() {
    Meter first(3);
    Meter second{4};
    use(first);
    use(Meter{5});
    // Meter bad = 3;   // blad: niejawna konwersja
    // Meter bad2 = {3}; // blad: copy-list-initialization
    // use(6);           // blad: niejawna konwersja argumentu

    Ready ready{true};
    if (ready) std::cout << "ready\n"; // kontekst wymaga bool
    bool flag = static_cast<bool>(ready); // jawna konwersja
    // bool bad = ready; // blad: operator jest explicit
    (void)second;
    (void)flag;
}
\end{lstlisting}

\texttt{explicit operator bool()} nadal działa w warunkach takich jak
\texttt{if}, \texttt{while} czy operand \texttt{!}: język dopuszcza tam
\emph{kontekstową konwersję do bool}. Nie pozwala natomiast na zwykłe
\texttt{bool flag = ready}. Słowo \texttt{explicit} nie zabrania
inicjalizacji bezpośredniej ani jawnego \texttt{static\_cast}.

Od C++20 można uzależnić tę własność od warunku obliczanego podczas
kompilacji: \texttt{explicit(condition)}. Prawdziwy warunek działa jak
\texttt{explicit}, a fałszywy jak jego brak:

\begin{lstlisting}[caption={Warunkowe explicit od C++20 (kompletny program).}]
template<bool AllowImplicit>
struct Token {
    explicit(!AllowImplicit) Token(int n) : value(n) {}
    int value;
};

int main() {
    Token<true> a = 3;  // dozwolona konwersja
    Token<false> b{3};  // jawna inicjalizacja
    // Token<false> c = 3; // blad: wymaga jawnej inicjalizacji
    return a.value + b.value == 6 ? 0 : 1;
}
\end{lstlisting}

\texttt{explicit} jest szczególnie przydatne dla typów jednostek,
identyfikatorów i uchwytów, w których przypadkowe utworzenie obiektu
z liczby lub napisu mogłoby ukryć błąd wywołania.

\section{Praktyczna organizacja plików: przykład całościowy}

Rozważmy licznik wspólny dla programu, stała kompilacyjna
oraz funkcję dostępną z kilku TU:

\begin{lstlisting}
// statistics.hpp
#pragma once
#include <cstddef>

inline constexpr std::size_t block_size = 1024;
extern unsigned total_jobs;

inline unsigned completed_jobs() {
    return total_jobs;
}
\end{lstlisting}

\begin{lstlisting}
// statistics.cpp
#include "statistics.hpp"

unsigned total_jobs = 0;
\end{lstlisting}

\begin{lstlisting}
// worker.cpp
#include "statistics.hpp"

void do_one_job() {
    ++total_jobs;
}
\end{lstlisting}

\begin{lstlisting}
// main.cpp
#include "statistics.hpp"
#include <iostream>

void do_one_job();

int main() {
    do_one_job();
    std::cout << block_size << " " << completed_jobs() << '\n';
}
\end{lstlisting}

Kompilacja przykładowa:
\begin{lstlisting}
g++ -std=c++20 -Wall -Wextra -pedantic \
    main.cpp worker.cpp statistics.cpp -o program
\end{lstlisting}

Mamy tu jedną definicję \texttt{total\_jobs} w \texttt{statistics.cpp},
jedną stałą \texttt{block\_size} udostępnioną przez
\texttt{inline constexpr}, a także legalną definicję funkcji
\texttt{completed\_jobs} w nagłówku. Przykład zakłada jeden wątek;
przy równoczesnych zapisach do licznika trzeba użyć synchronizacji,
którą omawia rozdział~\ref{chap:concurrency}.

\section{Najczęstsze błędy i reguły wyboru}

\subsection{Błędy}

\begin{enumerate}
    \item \textbf{Definicja zwykłej zmiennej w nagłówku:}
    zapis \texttt{int counter = 0;} powoduje wielokrotne definicje,
    jeśli nagłówek włączymy do kilku TU. Użyj \texttt{extern}
    i jednej definicji w \texttt{.cpp} albo \texttt{inline} (C++17).
    \item \textbf{\texttt{static} w nagłówku jako wspólny stan:}
    w rzeczywistości tworzy osobny stan w każdej TU.
    \item \textbf{Utożsamianie \texttt{const} z \texttt{constexpr}:}
    obiekt \texttt{const} może być inicjalizowany w runtime.
    \item \textbf{Utożsamianie \texttt{constexpr} z \texttt{consteval}:}
    zwykłą funkcję \texttt{constexpr} wolno wywoływać w runtime.
    \item \textbf{Utożsamianie \texttt{inline} z optymalizacją:}
    kompilator może wstawić lub nie wstawić ciała dowolnej funkcji,
    zgodnie z dozwolonymi optymalizacjami.
    \item \textbf{Utożsamianie \texttt{static} z synchronizacją:}
    bezpieczna inicjalizacja lokalnego statyka nie zabezpiecza
    późniejszych zapisów.
    \item \textbf{Używanie \texttt{volatile} zamiast \texttt{atomic}:}
    \texttt{volatile} nie eliminuje wyścigów danych.
    \item \textbf{Nadmierne korzystanie z globalnego mutowalnego stanu:}
    nawet poprawny pod względem ODR program może być trudny
    do testowania oraz podatny na problemy kolejności inicjalizacji.
\end{enumerate}

\subsection{Kolejność inicjalizacji obiektów globalnych}

Inicjalizacja dynamiczna obiektów o statycznym czasie przechowywania
w różnych TU może mieć nieoczywistą kolejność. Jest to źródło
\emph{static initialization order fiasco}. Jedną z praktycznych
technik jest przeniesienie zależności do lokalnego obiektu
statycznego, inicjalizowanego przy pierwszym wywołaniu:

\begin{lstlisting}
#include <string>

const std::string& application_name() {
    static const std::string name = "MyApplication";
    return name;
}
\end{lstlisting}

Gdy możliwe jest \emph{stałe inicjalizowanie}, preferuj proste
\texttt{constexpr} lub \texttt{constinit}. Pamiętaj jednak, że
\texttt{constinit} wyklucza inicjalizację dynamiczną, lecz
nie narzuca niemutowalności obiektu.

\subsection{Zalecenia w skrócie}

\begin{itemize}
    \item Potrzebujesz niemutowalnego dostępu? Użyj \texttt{const}.
    \item Potrzebujesz wyrażenia stałego? Użyj \texttt{constexpr}.
    \item Wywołanie musi być stałe? Rozważ \texttt{consteval}.
    \item Mutowalny obiekt ma wymagać inicjalizacji statycznej?
    Rozważ \texttt{constinit}.
    \item Chcesz ukryć szczegół implementacji w \texttt{.cpp}?
    Rozważ anonimową przestrzeń nazw.
    \item Potrzebujesz stanu zachowywanego między wywołaniami?
    Rozważ lokalną \texttt{static}, pamiętając o wątkach.
    \item Potrzebujesz jednego symbolu zdefiniowanego w \texttt{.cpp}?
    Deklaruj go w nagłówku przez \texttt{extern}.
    \item Potrzebujesz definicji w nagłówku współdzielonej przez TU?
    Użyj \texttt{inline}, jeśli jest odpowiednie; dla stałych
    kompilacyjnych \texttt{inline constexpr}.
    \item Potrzebujesz oddzielnego stanu wątku? Użyj \texttt{thread\_local}.
    \item Potrzebujesz komunikacji i synchronizacji między wątkami?
    Użyj atomików lub blokad, a nie \texttt{volatile}.
    \item Konstruktor lub operator konwersji mógłby uruchamiać
    niezamierzoną konwersję? Rozważ \texttt{explicit}.
\end{itemize}

\paragraph{Konkluzja.}
Najważniejsze pytanie nie brzmi ``które słowo kluczowe jest szybsze?'',
lecz: czy deklaracja ma określać \emph{czas życia}, \emph{łączność},
\emph{niemutowalność}, \emph{stałą ewaluację}, czy
\emph{dopuszczalność wielu definicji}? Dopiero wtedy wybór
specyfikatora jest jednoznaczny.
